% SpecMol paper — IMRAD skeleton.
% Build:
%   pdflatex main && bibtex main && pdflatex main && pdflatex main
% Tables under tables/ are auto-generated by `python paper/make_tables.py`.

\documentclass[11pt]{article}

\usepackage[utf8]{inputenc}
\usepackage[T1]{fontenc}
\usepackage{lmodern}
\usepackage{microtype}
\usepackage[margin=1in]{geometry}
\usepackage{booktabs}
\usepackage{amsmath, amssymb}
\usepackage{graphicx}
\usepackage{hyperref}
\usepackage{xcolor}
\usepackage{natbib}

\title{Learnable 3D Pair Representations for Spectral Molecular Property Prediction}
\author{Anonymous}
\date{}

\newcommand{\todo}[1]{\textcolor{red}{[TODO: #1]}}

\begin{document}
\maketitle

\begin{abstract}
\todo{Write abstract after Method and Results are stable.
Key claims:
(i) injecting Uni-Mol pair representations as learnable edge weights into a Chebyshev spectral GNN improves classification AUC on BBBP and BACE;
(ii) a static pair-to-edge gating (V2-T5) outperforms a dynamic pair-node update (T6) under matched compute;
(iii) fingerprint-only baselines remain a strong control and must be reported alongside graph variants.}
\end{abstract}

% ---------------------------------------------------------------------
\section{Introduction}
\label{sec:intro}

Molecular property prediction is a central problem in drug discovery and computational chemistry, and a long line of work has shown that accurate predictions hinge on representations that capture both the topology of chemical bonds and the geometry of three-dimensional conformers. Two-dimensional graph neural networks exploit bond-level structure but cannot directly access through-space interactions such as steric clashes, hydrogen bonds, or favorable stacking. Conversely, purely three-dimensional models depend on the quality of conformer generation and tend to be computationally expensive at the scale of high-throughput screening. A practical alternative is to start from a strong two-dimensional backbone and inject three-dimensional information only where it is most useful, but the design space for such an injection remains underexplored.

Self-supervised pretraining has produced increasingly capable molecular encoders along both modalities. On the two-dimensional side, contrastive frameworks such as MolCLR \citep{wang2022molclr} and graph transformers such as GROVER \citep{rong2020grover} learn transferable atom and bond representations from millions of unlabeled molecules. On the three-dimensional side, Uni-Mol \citep{zhou2023unimol} pretrains an SE(3)-equivariant encoder over conformer ensembles and produces, in addition to atomic embeddings, a dense pair representation $P \in \mathbb{R}^{N \times N \times d_p}$ that encodes geometric relationships between every atom pair. Multi-view pretraining approaches such as GraphMVP \citep{liu2022graphmvp} have further shown that aligning two- and three-dimensional views during pretraining benefits downstream prediction. The complementary question of how to consume a frozen three-dimensional pair representation at downstream time, without retraining the upstream encoder, has received considerably less attention.

Spectral graph neural networks offer a particularly clean entry point for such an injection. The recently introduced ChebNet II \citep{he2022chebnetii} parametrizes graph convolutions as Chebyshev polynomials of the symmetric normalized Laplacian and admits per-edge weights as a first-class input. This makes the edge weights of the Laplacian a natural locus at which to fold in three-dimensional information: an off-the-shelf pair representation can be projected to a scalar gate on each chemical bond, and the resulting weighted Laplacian then drives the spectral filtering. The challenge is that naively initializing such a gate to small or random values immediately perturbs the spectrum away from the unweighted baseline, destabilizing the contrastive pretraining objective and erasing the inductive bias of the two-dimensional backbone before any useful learning has taken place.

We address this with two designs that share the same backbone---a dual-path Chebyshev II encoder with low-pass and high-pass towers, fused with a fingerprint MLP and pretrained by an NT-Xent contrastive loss across four views---but differ in how three-dimensional pair information is fed into the spectral filter. The first, \textbf{V2-T5}, applies a small MLP followed by a symmetrized sigmoid with bias initialization $b{=}{+}5$ to the Uni-Mol pair representation, yielding an initial per-bond gate of approximately $0.993$. This near-identity initialization preserves the two-dimensional baseline at epoch zero and lets the gate gradually move away from one only when the gradient supports it. The second, \textbf{T6}, replaces this static gate with a bidirectional pair-node update that iteratively refines both representations at each layer, trading stability for additional capacity. We additionally report a fingerprint-only control that uses no graph encoder, because such a control is frequently omitted from molecular benchmarks but turns out to be highly competitive.

To compare the three variants fairly we adopt the Uni-Mol scaffold fold split with matched training budgets, seeds $\{9, 19, 29\}$, and reporting format across all runs. Under this protocol the static V2-T5 gating attains $0.862$ test AUC on BBBP, improving over the two-dimensional baseline by $2.6$ points and over the dynamic T6 variant by $2.2$ points; on FreeSolv it reduces test RMSE from $0.760$ to $0.675$. On BACE the fingerprint-only control reaches $0.846$ AUC and surpasses the two-dimensional baseline at $0.797$, underscoring that strong fingerprint baselines should not be omitted from molecular benchmarks. Taken together, these results suggest that a frozen three-dimensional pair representation is most useful when it is injected near the identity and allowed to remain close to it, rather than through an aggressive iterative refinement.

\paragraph{Contributions.}
\begin{itemize}
  \item A static pair-to-edge gating mechanism (V2-T5) that projects a frozen Uni-Mol pair representation onto the chemical-bond edge weights of a Chebyshev II spectral GNN, using a near-identity initialization that preserves contrastive pretraining stability.
  \item A dynamic bidirectional pair-node update variant (T6) evaluated under matched compute, together with an empirical analysis of why iterative refinement increases seed-level variance on small molecular benchmarks.
  \item A strict fair-comparison protocol on the Uni-Mol scaffold fold split with matched seeds, training budgets, and reporting format, including a fingerprint-only control that we find essential for honest assessment of graph-based molecular encoders.
\end{itemize}

% ---------------------------------------------------------------------
\section{Related Work}
\label{sec:related}

\paragraph{Spectral graph neural networks.}
Spectral GNNs trace back to the Chebyshev-polynomial formulation of \citet{defferrard2016chebnet}, which was simplified to a first-order propagation rule by \citet{kipf2017gcn} and is now the workhorse \texttt{GCN}. Subsequent work has either replaced spectral filters with attention-based aggregation \citep{velickovic2018gat} or strengthened expressiveness through injective neighborhood pooling \citep{xu2019gin}. A complementary line of research returned to the spectral viewpoint and asked which polynomial basis yields the most expressive filter: \citet{he2021bernnet} introduced Bernstein-basis spectral filters with non-negativity constraints, \citet{chien2021gprgnn} used a generalized PageRank polynomial, and \citet{he2022chebnetii} re-examined the Chebyshev basis and showed that learnable Chebyshev coefficients outperform earlier approximations across both homophilic and heterophilic graphs. Most recently, \citet{chen2024polygcl} adapted this polynomial-filter framework to graph contrastive learning. Our backbone follows this tradition: we adopt a dual-path Chebyshev~II encoder with explicit low- and high-pass towers.

\paragraph{Molecular graph neural networks.}
For molecular property prediction, \citet{gilmer2017mpnn} unified earlier graph models under the message-passing framework and introduced strong baselines for quantum-chemistry regression. \texttt{AttentiveFP} \citep{xiong2020attentivefp} popularized attention-based aggregation in drug discovery and remains a frequent baseline. These models operate on bond topology only and have no mechanism for through-space (non-bonded) interactions, which is precisely the gap that motivates injecting three-dimensional information into the spectral backbone.

\paragraph{Three-dimensional molecular representations.}
Three-dimensional models build inductive biases from atomic coordinates. \texttt{SchNet} \citep{schutt2017schnet} introduced continuous-filter convolutions over pairwise distances, \texttt{DimeNet} \citep{klicpera2020dimenet} added directional information through angle features, and \texttt{EGNN} \citep{satorras2021egnn} enforced exact $E(n)$ equivariance at the architecture level. More recently, \texttt{Uni-Mol} \citep{zhou2023unimol} pretrained an SE(3)-equivariant encoder over conformer ensembles and made available a dense pair representation $P \in \mathbb{R}^{N\times N\times d_p}$ that summarizes geometric relationships between every atom pair. Our work consumes this frozen Uni-Mol pair representation rather than training a new 3D encoder, so the contribution is orthogonal to the design of the 3D backbone itself.

\paragraph{Self-supervised pretraining on molecules.}
Self-supervised pretraining alleviates the small-data problem in molecular property prediction. \texttt{MolCLR} \citep{wang2022molclr} adopts contrastive learning on 2D molecular graphs, \texttt{GROVER} \citep{rong2020grover} adds graph-transformer self-supervision at scale, and \texttt{GraphMVP} \citep{liu2022graphmvp} aligns 2D and 3D views during pretraining. \texttt{GEM} \citep{fang2022gem} explicitly enriches molecular representations with geometry-aware self-supervised objectives. Our work is complementary to all of these: we do not propose a new SSL objective, but rather a downstream-time mechanism for plugging a frozen 3D pair representation into a 2D spectral backbone, evaluated under the matched-protocol regime advocated by the \texttt{MoleculeNet} benchmark \citep{wu2018moleculenet}.

% ---------------------------------------------------------------------
\section{Method}
\label{sec:method}

\subsection{Notation}
Let $G=(\mathcal{V},\mathcal{E})$ be a molecular graph with $N=|\mathcal{V}|$ heavy atoms and chemical bonds $\mathcal{E}\subseteq\mathcal{V}\times\mathcal{V}$. Each atom carries a feature vector $\mathbf{x}_i\in\mathbb{R}^{d_x}$ (atomic number, degree, hybridization, etc.); each bond an edge attribute $\mathbf{e}_{ij}\in\mathbb{R}^{d_e}$. We additionally consume a frozen pair representation $\mathbf{P}\in\mathbb{R}^{N\times N\times d_p}$ produced by a pretrained Uni-Mol encoder \citep{zhou2023unimol}, indexed over all atom pairs $(i,j)$. The chemical-bond adjacency $A\in\{0,1\}^{N\times N}$ is symmetric (each bond contributes both directed edges $(i,j)$ and $(j,i)$) and induces the symmetric normalized Laplacian $L_{\mathrm{sym}}=I-D^{-1/2}AD^{-1/2}$, where $D$ is the degree matrix.

\subsection{Backbone: Dual-Path Chebyshev Spectral GNN}
Our backbone, denoted \texttt{LH\_Direct}, processes $G$ along two parallel spectral paths and fuses them with a molecular fingerprint stream. Each spectral path is a Chebyshev~II encoder \citep{he2022chebnetii} that applies an order-$K$ polynomial filter to the normalized Laplacian:
\begin{equation}
    h^{(\ell)} = \sum_{k=0}^{K} \theta_k\, T_k(\tilde L)\, h^{(\ell-1)},
    \qquad \tilde L = 2L_{\mathrm{sym}}/\lambda_{\max} - I,
    \label{eq:chebnetii}
\end{equation}
where $T_k$ is the $k$-th Chebyshev polynomial and the coefficients $\{\theta_k\}_{k=0}^{K}$ are learnable. The two paths use independent coefficients tuned, by construction, towards low-pass and high-pass behavior respectively. The pooled graph representations $h_l, h_h\in\mathbb{R}^{d_h}$ are concatenated with a fingerprint embedding $h_{\mathrm{fp}}=\mathrm{MLP}(\mathrm{FP}(G))$, where $\mathrm{FP}$ concatenates PubChem, MACCS, and ErG fingerprints. The fused representation $z = [h_l; h_h; h_{\mathrm{fp}}]$ is the only object exposed to the downstream head.

\subsection{Variant V0: Two-Dimensional Baseline}
The V0 baseline runs \texttt{LH\_Direct} with $A$ unchanged, i.e.\ all bonds carry equal Laplacian weight. V0 ignores the Uni-Mol pair representation entirely and serves as the strict two-dimensional reference.

\subsection{Variant V2-T5: Static Pair-to-Edge Gating}
\label{sec:v2t5}
V2-T5 augments V0 with a learnable scalar gate on each chemical bond, computed from the frozen Uni-Mol pair representation. Concretely, for each bond $(i,j)\in\mathcal{E}$ we form
\begin{equation}
    g_{ij} = \sigma\!\left( \tfrac{1}{2}\bigl(\mathrm{MLP}_\phi(\mathbf{P}_{ij}) + \mathrm{MLP}_\phi(\mathbf{P}_{ji})\bigr) + b \right),
    \label{eq:v2t5gate}
\end{equation}
where $\mathrm{MLP}_\phi$ is a two-layer scalar-output network shared across edges, the symmetrization ensures $g_{ij}=g_{ji}$, and $\sigma$ is the logistic sigmoid. Crucially we initialize the bias term $b=+5$, which gives an initial gate of $\sigma(5)\approx 0.993$ across all bonds. This near-identity initialization preserves the V0 dynamics at epoch zero, so the contrastive pretraining objective starts from a near-optimal point and is free to move $g_{ij}$ away from $1$ only when gradient signal justifies it. The reweighted Laplacian $L_{\mathrm{sym}}(g)$ is then plugged into Eq.~(\ref{eq:chebnetii}) in place of $L_{\mathrm{sym}}$.

\subsection{Variant T6: Dynamic Pair-Node Update}
\label{sec:t6}
T6 trades the static gate of V2-T5 for an iterative, bidirectional pair-node update applied at each Chebyshev layer. Given current node features $h^{(\ell)}$ and pair features $\mathbf{P}^{(\ell)}$, T6 first produces a residual update $\Delta\mathbf{P}^{(\ell)}$ from a multi-head node-to-pair attention computation, then refreshes node features using the updated pair representation, and finally recomputes the edge gates via Eq.~(\ref{eq:v2t5gate}) on $\mathbf{P}^{(\ell+1)}$. The output projection of the pair-update block is zero-initialized so that $\mathbf{P}^{(0)}\equiv\mathbf{P}$, making T6 strictly more expressive than V2-T5 at initialization. The cost is one extra attention matmul per layer and an additional source of optimization variance, which we examine empirically in Section~\ref{sec:experiments}.

\subsection{Training Objective and Downstream Evaluation}
\label{sec:training}
\paragraph{Self-supervised pretraining.} All variants are pretrained with a four-view NT-Xent contrastive objective that pulls together low-pass, high-pass, mixed, and fingerprint embeddings of the same molecule and pushes apart embeddings of different molecules within a batch. We use temperature $\tau{=}0.5$, batch size $512$, and early-stopping based on training loss with patience $100$ epochs.

\paragraph{Downstream linear probe.} After pretraining, the encoder is frozen and a two-layer logistic regression head (with a single Dropout layer set to $p{=}0.2$ during training and disabled during evaluation) is trained on top of the fused representation $z$. For each pretraining seed we run three independent evaluation restarts with split seeds $\{9,19,29\}$ and report the test metric at the validation-best epoch. To prevent early-epoch noise from locking model selection on a lucky-spike checkpoint, we require both a minimum of $100$ probe epochs before any selection update and a strict improvement tolerance of $10^{-4}$.

% ---------------------------------------------------------------------
\section{Experiments}
\label{sec:experiments}

\subsection{Setup}
\paragraph{Datasets.} BBBP (binary classification, 2{,}039 molecules), BACE (binary classification, 1{,}513 molecules), FreeSolv (regression, 642 molecules). All splits follow the Uni-Mol scaffold fold protocol.
\paragraph{Seeds and budget.} \todo{Three training seeds $\{9,19,29\}$ with three evaluation restarts each; 1000 pretraining epochs, 2000 downstream epochs.}
\paragraph{Baselines.} V0 (2D-only), FP-only (fingerprint MLP head only, no graph encoder).

\subsection{Main Results}
% Auto-generated by paper/make_tables.py — do not edit by hand.
\begin{table}[t]
\centering
\caption{Main results across molecular property prediction benchmarks. Classification reports test AUC (higher is better); regression reports test RMSE (lower is better). All numbers are mean $\pm$ std across three training seeds with Uni-Mol fold splits.}
\label{tab:main}
\begin{tabular}{lccc}
\toprule
 & \multicolumn{2}{c}{Classification (AUC $\uparrow$)} & Regression (RMSE $\downarrow$) \\
\cmidrule(lr){2-3} \cmidrule(lr){4-4}
Variant & BBBP & BACE & FreeSolv \\
\midrule
V0 & 0.836 $\pm$ 0.035 & 0.797 $\pm$ 0.085 & -- \\
FP-only & -- & 0.846 $\pm$ 0.010 & -- \\
V2-T5 & 0.862 $\pm$ 0.025 & -- & -- \\
T6 & 0.840 $\pm$ 0.035 & -- & -- \\
\bottomrule
\end{tabular}
\end{table}

\todo{Discuss: V2-T5 best on BBBP (+2.6 over V0); FP-only strongest single-modality on BACE; BACE V2-T5/T6 numbers pending post-B4 re-run.}

\subsection{Ablations}
\todo{(a) bias init $b{=}0$ vs.\ $b{=}{+}5$; (b) Uni-Mol pair vs.\ random pair representation; (c) static (V2-T5) vs.\ dynamic (T6).}

% ---------------------------------------------------------------------
\section{Discussion}
\label{sec:discussion}

\todo{Why static beats dynamic on small datasets: dynamic update increases capacity but also variance (T6 seed-19 underperforms). Limitations: small benchmarks, only one 3D source (Uni-Mol), no large-scale HIV result yet.}

% ---------------------------------------------------------------------
\section{Conclusion}
\label{sec:conclusion}

\todo{Wrap up; future work: scale to HIV, multi-task ClinTox/Tox21, learnable rather than frozen 3D encoder.}

% ---------------------------------------------------------------------
\bibliographystyle{plainnat}
\bibliography{refs}

\end{document}
